% !TEX root =./ECE444_Practice_main.tex
%
% L5 practice set (Field Regions) -- the three terms of the infinitesimal-dipole
% field and where each wins, the two region boundaries at both size extremes
% (a 3 m dish and an 8 cm whip), the pi/8 phase-error criterion behind the range
% length, and why the pattern is still forming inside 2D^2/lambda.

\fullwidth{\textbf{Learning Objective 1.5: } I can identify and distinguish the
reactive near-field, radiating near-field, and far-field regions and calculate
the boundaries for a given antenna.
\\[4pt]\normalfont\small Show your work. A \textbf{Documentation} line at the end
is required (write \textbf{None} if you did not collaborate).}

\begin{questions}

%%%%%%%%%%%%%%%%%%%%%%%%%%%%%%%%%%%%%%%%%%%%%%%%%%%%%%%%%%%%%%%%%%%%%%%%%%%%
\question \textbf{Three terms in one equation.} The exact $E_\theta$ of an
infinitesimal dipole is a sum of three pieces whose relative sizes are set
entirely by $kr$:
\eq{ E_\theta \propto
     \underbrace{\dfrac{1}{r}}_{\text{radiation}} +
     \underbrace{\dfrac{1}{kr^{2}}}_{\text{induction}} +
     \underbrace{\dfrac{1}{k^{2}r^{3}}}_{\text{electrostatic}} }
Throughout this question the antenna operates at $f = 915\mhz$.

\begin{parts}

	\part The innermost region is called \emph{reactive}. Justify that word in one
	or two sentences, using what the complex Poynting vector does there.
		\begin{solution}[1.2in]
		Close in, $\vec{E}$ and $\vec{H}$ are $90\mydeg$ out of phase, so the radial
		Poynting vector $\tfrac12 E_\theta H_\phi^{*}$ is almost purely
		\emph{imaginary} --- and its imaginary part falls off as $1/r^{5}$, which is
		why it overwhelms everything near the antenna. Imaginary power is not
		transported: energy flows outward for a quarter cycle and all the way back the
		next. That is stored energy, the same reactive power you met at the terminals
		in Lesson 4, now living in the surrounding field.
		\end{solution}

	\begin{minipage}[t]{0.58\linewidth}
		\part Normalizing the radiation term to $1$, give the relative sizes of the
		three terms at $kr = 0.1$ and at $kr = 10$, and name the dominant term in each
		case.
		\begin{solution}[1.6in]
		Dividing through by the radiation term leaves powers of $1/kr$:
		\eq{
			\text{radiation} : \text{induction} : \text{electrostatic}
			&= 1 : \dfrac{1}{kr} : \dfrac{1}{(kr)^{2}} \\
			kr = 0.1: \quad &= 1 : 10 : 100 \\
			kr = 10: \quad &= 1 : 0.1 : 0.01
		}
		At $kr = 0.1$ the $1/r^{3}$ electrostatic term is $100\times$ the radiation
		term. At $kr = 10$ the $1/r$ radiation term is all that is left standing.
		\end{solution}
	\end{minipage}\hfill%
	\begin{minipage}[t]{0.37\linewidth}
		\raggedleft \vspace{12pt}
		$kr = 0.1$: \ansbox[1.5in]{$1{:}10{:}100$, $1/r^{3}$ wins}

		\vspace{4pt}
		$kr = 10$: \ansbox[1.5in]{$1{:}0.1{:}0.01$, $1/r$ wins}
	\end{minipage}

	\begin{minipage}[t]{0.58\linewidth}
		\part Find the distance at which the radiation and induction terms are equal.
		Give it in wavelengths and in centimetres at $915\mhz$.
		\begin{solution}[1.5in]
		The two are equal when $1/kr = 1$:
		\eq{
			kr &= 1 \quad\longrightarrow\quad r = \dfrac{1}{k} = \dfrac{\lambda}{2\pi}
				\approx 0.159\lambda \\
			\lambda &= \dfrac{c}{f} = \dfrac{3\ee{8}}{915\ee{6}} = 0.328\m \\
			r &= \dfrac{0.328}{2\pi} = 0.0522\m = 5.22\ \text{cm}
		}
		\end{solution}
	\end{minipage}\hfill%
	\begin{minipage}[t]{0.37\linewidth}
		\raggedleft \vspace{12pt}
		$r$ = \ansbox[1.5in]{$0.159\lambda = 5.22\ \text{cm}$}
	\end{minipage}

	\part A colleague measures this antenna's gain with a probe held $2\ \text{cm}$
	away. In two sentences, say what is wrong with the number he reports.
		\begin{solution}[1.3in]
		At $2\ \text{cm}$ he is well inside the $5.2\ \text{cm}$ crossover, where the
		stored $1/r^{2}$ and $1/r^{3}$ terms dominate --- most of what the probe picks
		up was never radiated and would have returned to the antenna anyway. Worse, the
		probe sits in the stored field and \emph{loads} the antenna, shifting its input
		impedance and its resonance, so the reading describes a two-body system rather
		than the antenna. Gain is defined in the far field; there is no valid gain
		number to be had at $2\ \text{cm}$.
		\end{solution}

\end{parts}

\newpage
%%%%%%%%%%%%%%%%%%%%%%%%%%%%%%%%%%%%%%%%%%%%%%%%%%%%%%%%%%%%%%%%%%%%%%%%%%%%
\question \textbf{An electrically large antenna.} An X-band ground-station
reflector has diameter $D = 3.0\m$ and operates at $f = 8\ghz$. The boundaries are
$r_{\text{react}} = 0.62\sqrt{D^{3}/\lambda}$ and $r_{\text{ff}} = 2D^{2}/\lambda$.

\begin{parts}

	\begin{minipage}[t]{0.58\linewidth}
		\part Find $\lambda$ and the electrical size $D/\lambda$. Is this antenna
		electrically large?
		\begin{solution}[1.2in]
		\eq{
			\lambda &= \dfrac{c}{f} = \dfrac{3\ee{8}}{8\ee{9}} = 0.0375\m \\
			\dfrac{D}{\lambda} &= \dfrac{3.0}{0.0375} = 80
		}
		Eighty wavelengths across --- emphatically large, so the $2D^{2}/\lambda$
		machinery is the right tool.
		\end{solution}
	\end{minipage}\hfill%
	\begin{minipage}[t]{0.37\linewidth}
		\raggedleft \vspace{12pt}
		$\lambda$ = \ansbox[1.4in]{$0.0375\m$}

		\vspace{4pt}
		$D/\lambda$ = \ansbox[1.4in]{$80$ (large)}
	\end{minipage}

	\begin{minipage}[t]{0.58\linewidth}
		\part Compute both region boundaries.
		\begin{solution}[1.7in]
		\eq{
			r_{\text{react}} &= 0.62\sqrt{\dfrac{(3.0)^{3}}{0.0375}}
				= 0.62\sqrt{720} \\
			&= 0.62(26.83) = 16.6\m \\
			r_{\text{ff}} &= \dfrac{2(3.0)^{2}}{0.0375} = \dfrac{18}{0.0375} = 480\m
		}
		\end{solution}
	\end{minipage}\hfill%
	\begin{minipage}[t]{0.37\linewidth}
		\raggedleft \vspace{12pt}
		$r_{\text{react}}$ = \ansbox[1.4in]{$\approx 16.6\m$}

		\vspace{4pt}
		$r_{\text{ff}}$ = \ansbox[1.4in]{$480\m$}
	\end{minipage}

	\begin{minipage}[t]{0.58\linewidth}
		\part The only outdoor range on base is $150\m$ long. Which region does that
		put the source in, and what fraction of the far-field distance is it?
		\begin{solution}[1.4in]
		\eq{ 16.6\m < 150\m < 480\m }
		so $150\m$ lands in the \textbf{radiating near-field} (Fresnel) region: the
		energy is genuinely leaving the antenna, but the pattern is still a function of
		range.
		\eq{ \dfrac{150}{480} = 0.31 }
		The range is at $31\%$ of the distance it needs to be.
		\end{solution}
	\end{minipage}\hfill%
	\begin{minipage}[t]{0.37\linewidth}
		\raggedleft \vspace{12pt}
		region = \ansbox[1.3in]{rad.\ near-field}

		\vspace{4pt}
		fraction = \ansbox[1.2in]{$0.31\ r_{\text{ff}}$}
	\end{minipage}

	\part You need this dish's sidelobe pattern and you cannot build a $480\m$
	range. Give two ways out and state what each costs.
		\begin{solution}[1.6in]
		\textbf{Near-field scanning}: probe the field on a surface a few wavelengths
		from the aperture and transform the measured near field into the far-field
		pattern mathematically. It fits indoors and is the standard answer for large
		apertures --- but it costs a precision positioner, careful probe correction, and
		hours of scan time, and it is only as good as your knowledge of the probe.
		\textbf{Compact range}: illuminate the dish with a shaped reflector that
		collimates the source into a locally plane wave over a quiet zone, buying you a
		flat wavefront in a $20\m$ chamber. It costs a large, expensive, precisely
		figured reflector, and the quiet zone limits the size of what you can test.
		(A defensible third: test at a lower frequency, since $r_{\text{ff}}$ scales as
		$1/\lambda$ --- but then you are not measuring the antenna you ship.)
		\end{solution}

\end{parts}

\newpage
%%%%%%%%%%%%%%%%%%%%%%%%%%%%%%%%%%%%%%%%%%%%%%%%%%%%%%%%%%%%%%%%%%%%%%%%%%%%
\question \textbf{An electrically small antenna.} A $915\mhz$ ISM whip on a
sensor node has largest dimension $D = 8\ \text{cm}$. Recall
$\lambda = 0.328\m$ at this frequency.

\begin{parts}

	\begin{minipage}[t]{0.58\linewidth}
		\part Compute $D/\lambda$, $2D^{2}/\lambda$, and $0.62\sqrt{D^{3}/\lambda}$.
		\begin{solution}[2.0in]
		\eq{
			\dfrac{D}{\lambda} &= \dfrac{0.08}{0.328} = 0.244 \\
			\dfrac{2D^{2}}{\lambda} &= \dfrac{2(0.08)^{2}}{0.328}
				= \dfrac{0.0128}{0.328} \\
			&= 0.039\m = 3.9\ \text{cm} \\
			0.62\sqrt{\dfrac{D^{3}}{\lambda}}
				&= 0.62\sqrt{\dfrac{5.12\ee{-4}}{0.328}} \\
			&= 0.62(0.0395) = 2.5\ \text{cm}
		}
		\end{solution}
	\end{minipage}\hfill%
	\begin{minipage}[t]{0.37\linewidth}
		\raggedleft \vspace{12pt}
		$D/\lambda$ = \ansbox[1.4in]{$0.244$}

		\vspace{4pt}
		$2D^{2}/\lambda$ = \ansbox[1.4in]{$3.9\ \text{cm}$}

		\vspace{4pt}
		$0.62\sqrt{D^{3}/\lambda}$ = \ansbox[1.4in]{$2.5\ \text{cm}$}
	\end{minipage}

	\part Your $2D^{2}/\lambda$ came out \emph{smaller than a wavelength}. Explain in
	two or three sentences why that number is not the far-field boundary here, and
	say what replaces it.
		\begin{solution}[1.5in]
		The $0.62\sqrt{D^{3}/\lambda}$ and $2D^{2}/\lambda$ formulas are aperture
		formulas: they assume the antenna is \emph{electrically large}
		($D > \lambda$), so that the spread in path length across the structure is what
		limits you. Here $D = 0.244\lambda$, so there is essentially no path-length
		spread to worry about and $2D^{2}/\lambda = 3.9\ \text{cm}$ is a statement about
		nothing. What limits you instead is the \textbf{reactive near field} of the
		element itself, which extends to about $\lambda/2\pi$ --- exactly the $kr = 1$
		crossover of Question 1. Beyond that, the radiation term is the only one left,
		so the far field effectively begins there.
		\end{solution}

	\begin{minipage}[t]{0.58\linewidth}
		\part Compute $\lambda/2\pi$ for this antenna, and state the distance at which
		you would actually stand to measure its pattern (use the $r \ge 5\lambda$
		practical rule).
		\begin{solution}[1.5in]
		\eq{
			\dfrac{\lambda}{2\pi} &= \dfrac{0.328}{2\pi} = 0.052\m = 5.2\ \text{cm} \\
			r_{\text{meas}} &\ge 5\lambda = 5(0.328) = 1.64\m
		}
		The $5\lambda$ rule buys margin: at $\lambda/2\pi$ the radiation term has only
		just taken over, while at $5\lambda$ the stored terms are down by a further
		factor of $30$ in field.
		\end{solution}
	\end{minipage}\hfill%
	\begin{minipage}[t]{0.37\linewidth}
		\raggedleft \vspace{12pt}
		$\lambda/2\pi$ = \ansbox[1.4in]{$5.2\ \text{cm}$}

		\vspace{4pt}
		$r_{\text{meas}}$ = \ansbox[1.4in]{$\ge 1.64\m$}

	\end{minipage}

	\part You are offered a $3\m$ anechoic chamber to characterize this whip. Is it
	long enough? Name the error source that actually limits the measurement, given
	your answer to (c).
		\begin{solution}[1.5in]
		Yes --- $3\m$ is nearly $10\lambda$, comfortably past the $1.64\m$ you need,
		and the aperture criterion never binds for an antenna this small. The binding
		error source is no longer \emph{distance}: it is everything else in the room.
		Chamber wall reflections, the positioner, and above all the common-mode current
		on the feed cable --- which on an antenna only $0.24\lambda$ long can radiate
		comparably to the whip itself --- set the accuracy floor. Choke the cable and
		check the quiet-zone reflectivity; do not spend money on a longer chamber.
		\end{solution}

\end{parts}

\newpage
%%%%%%%%%%%%%%%%%%%%%%%%%%%%%%%%%%%%%%%%%%%%%%%%%%%%%%%%%%%%%%%%%%%%%%%%%%%%
\question \textbf{Why ranges are long: the $\pi/8$ criterion.} A source a
distance $r$ away is farther from the \emph{edge} of an aperture of size $D$ than
from its centre, by a path difference $\Delta \approx D^{2}/(8r)$. The agreed
tolerance is $\Delta \le \lambda/16$. Take the lesson's reflector:
$D = 1.2\m$ at $f = 10\ghz$, so $\lambda = 0.03\m$.

\begin{parts}

	\begin{minipage}[t]{0.58\linewidth}
		\part At $r = 30\m$, find the path difference $\Delta$ and the phase error it
		produces across the aperture, in degrees.
		\begin{solution}[1.6in]
		\eq{
			\Delta &= \dfrac{D^{2}}{8r} = \dfrac{(1.2)^{2}}{8(30)}
				= \dfrac{1.44}{240} = 6.0\ \text{mm} \\
			\Delta\phi &= 2\pi \dfrac{\Delta}{\lambda}
				= 2\pi \lp \dfrac{0.006}{0.03} \rp = 1.257 \text{ rad} \\
			&= 72\mydeg
		}
		\end{solution}
	\end{minipage}\hfill%
	\begin{minipage}[t]{0.37\linewidth}
		\raggedleft \vspace{12pt}
		$\Delta$ = \ansbox[1.4in]{$6.0\ \text{mm}$}

		\vspace{4pt}
		$\Delta\phi$ = \ansbox[1.4in]{$72\mydeg$}
	\end{minipage}

	\begin{minipage}[t]{0.58\linewidth}
		\part Impose $\Delta \le \lambda/16$ and solve for the minimum range. Show
		that the $\pi/8$ tolerance is what produces the familiar $2D^{2}/\lambda$.
		\begin{solution}[1.7in]
		\eq{
			\dfrac{D^{2}}{8r} &\le \dfrac{\lambda}{16}
				\quad\longrightarrow\quad r \ge \dfrac{16D^{2}}{8\lambda}
				= \dfrac{2D^{2}}{\lambda} \\
			r_{\text{ff}} &= \dfrac{2(1.2)^{2}}{0.03} = 96\m
		}
		A path difference of $\lambda/16$ is a phase error of
		$2\pi/16 = \pi/8 = 22.5\mydeg$ --- so ``far field'' and ``no more than
		$22.5\mydeg$ of quadratic phase across the aperture'' are the same statement.
		\end{solution}
	\end{minipage}\hfill%
	\begin{minipage}[t]{0.37\linewidth}
		\raggedleft \vspace{12pt}
		$r_{\text{ff}}$ = \ansbox[1.4in]{$96\m$}
	\end{minipage}

	\part Your chamber is the $30\m$ one from part (a) and the program will not fund
	a $96\m$ outdoor range. State what the $72\mydeg$ of residual phase error does to
	the measured pattern, and decide whether you would accept the $30\m$ data for a
	\emph{gain} check, for a \emph{beamwidth} check, and for a \emph{sidelobe} check.
		\begin{solution}[1.8in]
		Quadratic phase across the aperture behaves like a defocus: it spoils the
		coherent addition at the edges, so the measured main lobe broadens slightly and
		loses a few tenths of a dB of peak, the nulls \emph{fill in}, and the near-in
		sidelobes \emph{rise} --- the last effect being much the largest.
		\textbf{Gain}: marginal, and it will read low; usable only with a stated
		correction. \textbf{Beamwidth}: acceptable, since the main lobe is the least
		sensitive feature. \textbf{Sidelobes}: reject --- at $72\mydeg$ the sidelobe
		structure is an artifact of the range, and a $-25\db$ requirement cannot be
		verified this way. If sidelobes are the deliverable, pay for near-field
		scanning.
		\end{solution}

	\part Why is $22.5\mydeg$ ``good enough''? Answer in one or two sentences, saying
	what would change if the tolerance were tightened to $\pi/16$.
		\begin{solution}[1.3in]
		It is a deliberate engineering compromise: at $22.5\mydeg$ the residual pattern
		error is a few tenths of a dB in gain and a fraction of a dB in the first
		sidelobes --- small compared with the other uncertainties on a real range --- so
		nothing is gained by going further. Halving the tolerance to $\pi/16$ would
		\emph{double} the required range length, because $r$ goes inversely with the
		allowed phase error, and $96\m$ would become $192\m$ for an improvement no one
		could measure.
		\end{solution}

\end{parts}

\newpage
%%%%%%%%%%%%%%%%%%%%%%%%%%%%%%%%%%%%%%%%%%%%%%%%%%%%%%%%%%%%%%%%%%%%%%%%%%%%
\question \textbf{Two antennas on one airframe.} A $3\ghz$ ($\lambda = 0.1\m$)
installation carries a half-wave dipole ($D = \lambda/2 = 0.05\m$) and a $2.0\m$
reflector, mounted $10\m$ apart.

\begin{parts}

	\part Inside $2D^{2}/\lambda$ the field is genuinely radiating --- energy is
	leaving and never coming back --- and yet the \emph{pattern} still changes as you
	walk outward. Explain why, in words.
		\begin{solution}[1.7in]
		The pattern is a superposition: every patch of the aperture sends a
		contribution to your observation point, and what you measure is those
		contributions adding with their relative phases. Close in, the patches are at
		meaningfully \emph{different} distances from you, so those relative phases
		depend on where you are standing --- move outward and the interference re-forms
		into a different angular pattern. As $r$ grows the quadratic phase across the
		aperture, $\propto D^{2}/r$, shrinks; past $2D^{2}/\lambda$ it is below
		$22.5\mydeg$ and the rays are parallel enough that the relative phases stop
		depending on $r$ altogether. From there the pattern is frozen and only the
		amplitude keeps falling as $1/r$.
		\end{solution}

	\begin{minipage}[t]{0.58\linewidth}
		\part Compute $2D^{2}/\lambda$ for each antenna.
		\begin{solution}[1.5in]
		\eq{
			\text{dipole:} \quad \dfrac{2(0.05)^{2}}{0.1} &= \dfrac{0.005}{0.1}
				= 0.05\m \\
			\text{dish:} \quad \dfrac{2(2.0)^{2}}{0.1} &= \dfrac{8}{0.1} = 80\m
		}
		A factor of $1600$ between two antennas sharing a frequency --- the far-field
		distance is set by \emph{size}, not by wavelength alone.
		\end{solution}
	\end{minipage}\hfill%
	\begin{minipage}[t]{0.37\linewidth}
		\raggedleft \vspace{12pt}
		dipole = \ansbox[1.4in]{$0.05\m$}

		\vspace{4pt}
		dish = \ansbox[1.4in]{$80\m$}
	\end{minipage}

	\begin{minipage}[t]{0.58\linewidth}
		\part At the $10\m$ separation, which region is each antenna in with respect
		to the other?
		\begin{solution}[1.7in]
		The dish sits at $10\m$ from the dipole, and $10\m \gg 0.05\m$ --- it is deep
		in the dipole's \textbf{far field}. The dipole sits at $10\m$ from the dish,
		well short of $80\m$. Check the inner boundary:
		\eq{ 0.62\sqrt{\dfrac{(2.0)^{3}}{0.1}} = 0.62\sqrt{80} = 5.5\m }
		Since $5.5\m < 10\m < 80\m$, the dipole is in the dish's \textbf{radiating
		near-field}. The relationship is not symmetric.
		\end{solution}
	\end{minipage}\hfill%
	\begin{minipage}[t]{0.37\linewidth}
		\raggedleft \vspace{12pt}
		dish: \ansbox[1.4in]{dipole's far field}

		\vspace{4pt}
		dipole: \ansbox[1.4in]{dish's rad.\ near fld}
	\end{minipage}

	\part The EMC engineer wants to predict the coupling between the two with Friis.
	Is that legitimate here? Say what you would do instead, and what it costs.
		\begin{solution}[1.6in]
		No. Friis assumes each antenna sees the other as a point source radiating its
		far-field pattern --- it needs \emph{both} ends in the far field, and the dipole
		is $8\times$ too close to the dish for that. Used anyway, Friis would quote the
		dish's boresight gain at a range where that gain has not formed, and the answer
		could be wrong by many dB in either direction.
		Do it properly with a full-wave simulation of the installed geometry (method of
		moments or FDTD over the airframe), or measure the coupling directly with a
		network analyzer on the real structure. Both cost far more than an equation:
		the simulation needs an accurate airframe model and hours of solve time, and the
		measurement needs the actual aircraft.
		\end{solution}

\end{parts}

\vspace{1.5em}
\noindent\textbf{Documentation:} \rule[-2pt]{0.6\linewidth}{0.4pt}

\end{questions}
